以下结果推广了 Gårding 不等式.

\begin{Theorem}
\label{thm:agmon-boundary-constrained-coercivity}
设 $L$, $B_j$ 和 $a$ 如上. 假设在 $\partial\Omega$ 的每一点 Agmon 条件均成立. 则存在常数 $c_1$ 和 $c_2$, 使 \eqref{eq:agmon-boundary-constrained-coercive-estimate} 成立.
\end{Theorem}

现在讨论如何将 \eqref{eq:agmon-shifted-constrained-variational-equation} 解释为椭圆边值问题. 为此, 首先需要一个正则性陈述.

\begin{Theorem}
\label{thm:agmon-weak-variational-solution-regularity}
假设 $\Omega$ 以及 $L$ 和 $B_j$ 的系数足够光滑. 再假设 $f\in L^2(\Omega)$. 则 \eqref{eq:agmon-shifted-constrained-variational-equation} 的解 $u$ 属于 $H^{2m}(\Omega)$.
\end{Theorem}

接着需要一个 Green 公式.

\begin{Theorem}
\label{thm:agmon-constrained-form-green-boundary-identity}
设 $L$ 和 $a$ 如上. 设 $B_i(\bm{x},D)$, $i=1,\ldots,m$, 为 $m$ 阶 Dirichlet 系统. 假设 $\Omega$ 以及所涉及算子的系数足够光滑. 则存在阶数为 $2m-1-\operatorname{ord}B_i$ 的正规边值算子 $C_i$, 使得对所有 $u,v\in H^{2m}(\Omega)$ 有
\begin{equation}\label{eq:agmon-constrained-form-green-boundary-identity}
a(u,v)=\int_\Omega(Lu)v\,\mathrm{d}\bm{x}-\sum_{i=1}^m\int_{\partial\Omega}(C_iu)(B_iv)\,\mathrm{d}S.
\end{equation}

\end{Theorem}

证明与 Theorem~\ref{thm:elliptic-green-formula-adjoint-boundary-operators} 的证明完全类似. 对 $u\in H^{2m}(\Omega)$ 和 $f\in L^2(\Omega)$, 方程 \eqref{eq:agmon-shifted-constrained-variational-equation} 现在成为
\begin{equation}\label{eq:agmon-shifted-form-natural-boundary-reduction}
\int_\Omega(Lu+\lambda u)v\,\mathrm{d}\bm{x}-\sum_{j=p+1}^m\int_{\partial\Omega}(C_ju)(B_jv)\,\mathrm{d}S=\int_\Omega fv\,\mathrm{d}\bm{x}.
\end{equation}
这表明 \eqref{eq:agmon-shifted-constrained-variational-equation} 是椭圆边值问题
\begin{equation}\label{eq:agmon-essential-natural-boundary-value-problem}
Lu+\lambda u=f,\ B_ju=0,\ j=1,\ldots,p,\ C_ju=0,\ j=p+1,\ldots,m.
\end{equation}
的弱形式.

第一组边界条件称为本质边界条件; 在问题的弱表述中, 直接将它们施加于 $u$. 第二组边界条件称为 ``自然\textquotesingle\textquotesingle\ 边界条件; 它们不显式施加, 而是像 Section~\ref{subsec:neumann-problem} 中的 Neumann 条件那样由分部积分产生.
